\chapter{開始使用 \headinglatin{KSSOLV Toolbox}}\label{ch:introduction}
\chapterlead{KSSOLV Toolbox 把原子結構、計算步驟和結果組織在同一個工作區中。可以從一個分子或晶體開始，檢查它的幾何形狀，選擇計算方法，再沿著工作流得到總能量、電子能級或其他結果。本章說明本書的閱讀方法，以及完成一次計算時需要建立的幾個基本認識。}

\section{從結構走向結果}
對一個材料體系提出問題，往往先於開啟軟體。例如，改變晶格常數會怎樣影響總能量？一個缺陷會在晶體中引入怎樣的區域性環境？分子的初始鍵長經過最佳化後會如何變化？KSSOLV Toolbox 為這些研究提供結構準備和電子結構計算工具。它使用平面波基組和贗勢描述電子體系，由 KSSOLV 完成 Kohn--Sham 方程的數值求解。

開始練習時，不必先掌握全部理論。應先認清三件事：計算物件是什麼，軟體將執行哪些步驟，最終的量回答了什麼問題。以氫分子為例，結構給出兩個氫原子的位置；自洽場計算在這些位置不變時求電子基態；結構最佳化進一步調整原子座標，尋找能量較低的幾何形狀。這些步驟之間的區別，比一次輸入多少引數更值得先弄清。

\subsection{完成一次計算的基本步驟}
\begin{steps}
\item \term{準備結構。}匯入檔案，或使用建模工具建立結構。檢查元素、原子數、晶格及週期性。
\item \term{建立工作流。}從模板開始，選擇結構輸入、計算設定、求解和繪圖等節點。
\item \term{設定計算條件。}確認贗勢、截斷能、$k$ 點、收斂閾值和需要的結果。
\item \term{執行並檢查。}檢視輸出訊息、迭代過程和結束狀態，判斷求解是否完成。
\item \term{分析與儲存。}閱讀數值與圖形，記錄輸入條件，將工作流和結果儲存在專案中。
\end{steps}

這個順序也適用於稍複雜的任務。表面吸附需要先建立有真空層的表面模型；能帶計算需要先獲得自洽電子密度；遠端執行則在提交前選好計算位置。後續章節會把這些條件放回具體例項中說明。

\section{選擇使用方式}
KSSOLV Toolbox 提供獨立應用程式、MATLAB\regmark{} 工具箱和\ksonline{}三種使用方式。它們共享結構、專案和工作流的主要概念，選擇時主要考慮現有軟體環境、檔案儲存位置以及計算在哪臺計算機上執行。

\begin{table}[!htb]
\centering
\caption{三種使用方式的執行條件與主要用途。}
\label{tab:usage-modes}
\begin{tabularx}{\linewidth}{>{\RaggedRight\arraybackslash}p{28mm}>{\RaggedRight\arraybackslash}X>{\RaggedRight\arraybackslash}X}
\toprule
使用方式 & 執行條件 & 適合的使用場景 \\
\midrule
獨立應用程式 & 無須安裝完整 MATLAB\regmark{}，也無須購買或持有 MATLAB\regmark{} 許可證；需要與應用發行包匹配的 MATLAB\regmark{} Runtime。 & 從本機系統直接啟動，在本機準備結構、執行計算並儲存專案。 \\
工具箱版本 & 已有受支援版本的 MATLAB\regmark{} 及其有效許可證；要求 MATLAB\regmark{} R2025a 或更新版本。 & 在圖形工作區中操作，同時結合 MATLAB\regmark{} 變數、指令碼及其他已授權的 MATLAB\regmark{} 產品。 \\
\ksonline{} & 透過瀏覽器使用，無須安裝客戶端；工作空間與計算會話執行在伺服器上。 & 使用統一配置的伺服器環境，上傳輸入、管理雲端專案並下載結果。 \\
\bottomrule
\end{tabularx}
\end{table}

\subsection{獨立應用程式}
獨立應用程式安裝在本機，由 MATLAB\regmark{} Runtime 提供執行環境。專案、結構與匯出檔案通常儲存在本機磁碟，檔案選擇視窗使用作業系統的桌面環境。完成 Runtime 準備後，可以從系統應用選單或快捷方式啟動，無須先開啟 MATLAB\regmark{}。

\subsection{\headinglatin{MATLAB}\regmark{} 工具箱}
工具箱安裝在受支援的 MATLAB\regmark{} 中，既可使用圖形介面，也可與 MATLAB\regmark{} 工作區和指令碼配合。結構匯入、建模、工作流配置與結果檢視仍從同一套工作區進入。本書會在涉及 MATLAB\regmark{} 變數、額外產品或遠端作業時說明所需條件，隨書指令碼則提供復算輸入與參考結果。

\subsection{\headinglatin{\ksonline}}
\ksonline{} 在瀏覽器中提供工作空間，瀏覽器負責顯示和互動，伺服器會話負責執行計算。瀚海量子\regmark{}官方線上版即將推出，並支援叢集私有化部署。

線上工作目錄與本機下載目錄是兩個位置。上傳把本機輸入送入伺服器，儲存專案則把當前工作寫入雲端目錄；需要本機副本時，再下載專案或匯出結果。登入、檔案傳輸和會話操作見第\ref{ch:online}章。

\section{理解專案結構}
\term{專案}是一次工作的容器，副檔名為 \file{.ks}。一個專案可以包含多個結構、多個工作流及其結果。可將它理解為一本實驗記錄：既收錄研究物件，也保留計算的組織方式。

\term{結構}說明原子是什麼、位於何處，以及週期結構的晶胞如何定義。三維檢視便於檢查幾何，但螢幕中的球體大小、背景顏色和顯示重複範圍不是電子結構計算引數。

\term{工作流}說明步驟如何相接。節點代表某項任務，連線表示資料的傳遞。把結構與 SCF 計算相連，可以讓求解器讀取它；把計算結果接到收斂圖節點，可以檢視迭代過程。工作流本身也可儲存成 \file{.wf} 檔案，便於在其他專案中重用。

\term{資料與圖形}儲存在專案瀏覽器的結果分支中。資料可以用於後續計算，圖形便於觀察變化。僅儲存一張圖片，通常不足以重新分析結果；應同時保留專案與必要的原始資料。

\begin{notebox}{先區分“看見的結構”和“參與計算的結構”}
將晶體的顯示重複設為 $2\times2\times2$，可以幫助觀察相鄰晶胞，但不會自動建立八倍原子數的計算模型。若要研究真實超胞中的空位或摻雜，應使用\ui{構建超胞}。這一區別將在建模教程中實際演示。
\end{notebox}

\section{本書的閱讀路線}
\subsection{第一次完成計算}
先閱讀安裝、線上版和工作區章節，接著完成第\ref{ch:first}章的氫分子計算。練習涵蓋結構匯入、開啟模板、設定基本引數、執行求解與儲存結果。隨後閱讀計算基礎，可以把已經見過的引數與物理含義聯絡起來。

\subsection{準備材料模型}
如果已經熟悉計算，但需要構建自己的材料，可從第\ref{ch:modeling}章進入。矽晶體、苯環和乙醇分別展示週期結構與分子的日常操作；第\ref{ch:advanced}章把這些方法延伸到表面、異質結、扭轉雙層、奈米結構和軟物質，第\ref{ch:modeling-tools}章則說明具體工具的引數與使用方法。每次改變結構後，先核對原子數、晶胞和區域性距離，再把模型交給計算工作流。

\subsection{整理與展示結果}
第\ref{ch:volume}章介紹體資料，第\ref{ch:results}章介紹數值和圖形結果。閱讀時應把結果與輸入條件放在一起：總能量屬於哪一個結構，能帶採用哪條路徑，等值面來自哪個資料集。這樣得到的圖才具有清楚的含義，也便於在研究報告中說明計算依據。

\section{本書約定}
\ui{主頁 → 結構 → 從檔案匯入結構}表示依次選擇頁面、按鈕和選單項。若選單仍顯示英文，本書保留英文名稱。\file{silicon.cif}等名稱表示檔案；鍵盤上的按鍵寫成 \ui{Enter}、\ui{Esc} 或 \ui{Ctrl/Cmd+Z}。後者表示 Windows/Linux 常用 Ctrl，macOS 常用 Command。

本書隨附的 \file{examples} 目錄包含教程檔案。練習前將它複製到自己的工作目錄，例如 \file{KSSOLV-Study}。建議為每一例建立獨立子目錄，使輸入結構、專案與匯出圖形儲存在同處。檔名可以體現物件和用途，例如 \file{H2-SCF.ks}、\file{Si-vacancy.ks}；這樣在數週後重新開啟時，也容易認出內容。

教程中的引數用於演示操作與解釋結果。研究自己的體系時，應圍繞目標物理量進行收斂檢查。相關章節介紹如何逐項調整引數、比較結果，判斷精度是否滿足需要。

\section{從一個清楚的問題開始}
軟體可以很快生成複雜的結構和工作流，而可靠的練習通常從簡單物件開始。先完成一個已知的小體系，再更換材料；先確認自洽計算，再增加能帶或態密度步驟；先儲存一個清楚的專案，再嘗試批次操作。每一步都應能用一句話說明目的。

閱讀本書時，可以把關鍵操作的檢查結果寫進自己的實驗記錄。記錄不必很長：原子數、晶胞、主要引數、是否收斂，以及一兩條觀察就足夠。隨著練習深入，這些簡短記錄會幫助把不同階段的模型與結果聯絡起來。


需要查詢具體任務時，可用第\ref{ch:workflows}章的工作流說明與附錄任務速查，從所需結果查到輸入條件和相應教程。
